\section{Virtualização}

Para computação, virtualização é a inserção de uma camada adicional entre um \gls{mediador-de-hardware} e o \gls{hardware} físico. Quando um componente é virtualizado,
esse recurso específico deixa de ser utilizado de forma direta e passa a ser apresentado por meio de uma abstração lógica controlada por essa camada intermediária.
Isso permite, por exemplo, que múltiplas aplicações utilizem simultaneamente o mesmo recurso, pois a camada de virtualização fica responsável por arbitrar o acesso
e coordenar seu compartilhamento evitando conflitos.

\section{Hipervisor}
Enquanto \Glspl{so} tradicionalmente atuam como supervisores (\autoref{fig:niveis-privilegio}), um hipervisor introduz uma camada adicional de abstração.
Ele opera um nível de privilégio abaixo dos \glspl{so} convidados, bloqueando o acesso direto deles ao hardware e virtualizando os recursos físicos para
orquestrar múltiplos \glspl{so} de forma isolada.

Para que essa estrutura funcione, o hipervisor utiliza mecanismos de interceptação. Como o \gls{so} convidado opera acreditando deter o privilégio máximo
sobre a máquina, ele tenta executar instruções privilegiadas diretamente. No entanto, como ele não possui esse controle essas ações geram
uma exceções que são capturadas pelo hipervisor. Este, por sua vez, traduz ou emula a operação, garantindo que o acesso aos recursos físicos ocorra
de forma segura e isolada.

Fundamentalmente, existem dois tipos de hipervisors: 1 e 2.

\subsection{Tipo 1}

Hipervisores de tipo 1, também conhecidos como bare-metal, são aqueles que executam diretamente sobre o hardware físico, sem a necessidade de um
sistema operacional hospedeiro intermediário. Eles atuam como \gls{so} e como camada de virtualização apresentando maior desempenho e menor latência
em comparação aos hipervisores do Tipo 2. Além disso, proporcionam maior controle sobre os recursos do sistema e melhor isolamento entre as máquinas virtuais,
tornando-os especialmente adequados para ambientes críticos.

\begin{figure}[htbp]
    \centering
    \begin{tikzpicture}[
        pequeno/.style={
            draw,
            thick,
            fill=blue!5,
            minimum width=4cm,
            minimum height=1.2cm,
            outer sep=0pt
        }
    ]
        \node[pequeno] (apps-so1-1) {Aplicação 1};
        \node[pequeno, right=0.2cm of apps-so1-1] (apps-so1-2) {Aplicação 2};
        \node[pequeno, right=0.2cm of apps-so1-2] (apps-so2-1) {Aplicação 1};

        \node[fit=(apps-so1-1)(apps-so1-2), inner sep=0pt] (apps-so1) {};
        \node[fit=(apps-so2-1), inner sep=0pt] (apps-so2) {};

        \node[
            draw,
            thick,
            fill=blue!10,
            minimum width=8.2cm,
            minimum height=1.5cm,
            below=0.2cm of apps-so1
		] (SO1) {\Gls{so} 1};

		\node[
            draw,
            thick,
            fill=blue!10,
            minimum width=4cm,
            minimum height=1.5cm,
            below=0.2cm of apps-so2
		] (SO2) {\Gls{so} 2};

		\node[fit=(SO2)(SO1), inner sep=0pt] (SOs) {};

		\node[
            draw,
            thick,
            fill=blue!15,
            minimum width=12.4cm,
            minimum height=1.5cm,
            below=0.2cm of SOs
		] (HOS) {Hipervisor};

		\node[
            draw,
            thick,
            fill=blue!20,
            minimum width=12.4cm,
            minimum height=1.5cm,
            below=0.2cm of HOS
		] () {\Gls{hardware}};

    \end{tikzpicture}
    \caption{Hipervisor Tipo 1.}
    \label{fig:hypervisor-type-1}
\end{figure}

\subsection{Tipo 2}

Diferentemente da solução nativa, um hipervisor de Tipo 2 é executado como um processo dentro de um \gls{so} hospedeiro.
Apesar de garantir o isolamento e a virtualização para diversos \glspl{so} convidados, sua arquitetura impõe mais uma camada intermediária de
software que restringe a eficiência no acesso aos componentes físicos, gerando latência adicional e menor desempenho em operações críticas.

\begin{figure}[htbp]
    \centering
    \begin{tikzpicture}[
        pequeno/.style={
            draw,
            thick,
            fill=blue!5,
            minimum width=4cm,
            minimum height=1.2cm,
            outer sep=0pt
        }
    ]
        \node[pequeno] (apps-so1-1) {Aplicação 1};
        \node[pequeno, right=0.2cm of apps-so1-1] (apps-so1-2) {Aplicação 2};
        \node[pequeno, right=0.2cm of apps-so1-2] (apps-so2-1) {Aplicação 1};

        \node[fit=(apps-so1-1)(apps-so1-2), inner sep=0pt] (apps-so1) {};
        \node[fit=(apps-so2-1), inner sep=0pt] (apps-so2) {};

        \node[
            draw,
            thick,
            fill=blue!10,
            minimum width=8.2cm,
            minimum height=1.5cm,
            below=0.2cm of apps-so1
		] (SO1) {\Gls{so} 1};

		\node[
            draw,
            thick,
            fill=blue!10,
            minimum width=4cm,
            minimum height=1.5cm,
            below=0.2cm of apps-so2
		] (SO2) {\Gls{so} 2};

		\node[fit=(SO2)(SO1), inner sep=0pt] (SOs) {};

		\node[
            draw,
            thick,
            fill=blue!15,
            minimum width=12.4cm,
            minimum height=1.5cm,
            below=0.2cm of SOs
		] (HOS) {Hipervisor};

		\node[
            draw,
            thick,
            fill=blue!20,
            minimum width=12.4cm,
            minimum height=1.5cm,
            below=0.2cm of HOS
		] (OS) {\Gls{so}};

		\node[
            draw,
            thick,
            fill=blue!25,
            minimum width=12.4cm,
            minimum height=1.5cm,
            below=0.2cm of OS
		] () {\Gls{hardware}};

    \end{tikzpicture}
    \caption{Hipervisor Tipo 2.}
    \label{fig:hypervisor-type-2}
\end{figure}




\section{Interceptação e Emulação - Trap and Emulate}



